\newpage
\section{Revisão bibliográfica}
\label{sc:revisao}

A literatura recente sobre Geração Aumentada de Recuperação (RAG) organiza-se em torno de um problema central: como integrar conhecimento externo a Modelos de Linguagem de Grande Porte (LLMs) de forma simultaneamente precisa e computacionalmente sustentável. Para situar a contribuição deste trabalho, esta revisão estrutura-se em cinco eixos temáticos: (i) a motivação da arquitetura RAG frente às limitações dos LLMs; (ii) os métodos de manipulação de contexto e seus compromissos de eficiência; (iii) as estratégias avançadas de recuperação; (iv) a infraestrutura de bancos de dados vetoriais; e (v) as aplicações de domínio específico. Ao fim, discutem-se as lacunas que justificam a investigação proposta.

\subsection{Limitações dos LLMs e a consolidação do RAG}

Os LLMs, apesar de seu desempenho notável, são propensos a alucinações e à desatualização do conhecimento paramétrico. O levantamento de Andriopoulos e Pouwelse~\cite{andriopoulos2023augmenting} sistematiza as estratégias de aumento de conhecimento como principal caminho para prevenir alucinações, posicionando o RAG como uma alternativa não-paramétrica de baixo custo de atualização frente ao retreinamento contínuo dos modelos. Essa fragilidade também se manifesta na própria engenharia dos sistemas: Barnett et al.~\cite{barnett2024seven}, a partir de três estudos de caso em domínios distintos (pesquisa, educação e biomedicina), catalogam sete pontos de falha recorrentes na construção de pipelines RAG, concluindo que a robustez desses sistemas é construída empiricamente durante a operação e não estabelecida em projeto. Em conjunto, esses trabalhos delimitam o problema que motiva a área, mas tratam predominantemente da acurácia e da confiabilidade, sem quantificar de forma sistemática o custo de hardware envolvido.

\subsection{Métodos de manipulação de contexto e o compromisso de eficiência}

O núcleo metodológico deste estudo apoia-se na caracterização dos métodos de injeção de contexto. Topsakal e Akinci~\cite{topsakal2023langchain} descrevem, em nível arquitetural, o funcionamento dos quatro fluxos fundamentais disponíveis em frameworks como o LangChain --- Stuff, Map Reduce, Refine e Map Re-rank ---, evidenciando trade-offs estruturais como o gargalo de paralelização do Refine e o elevado volume de chamadas do Map Reduce. Trata-se, contudo, de um trabalho descritivo, sem medições empíricas. Essa lacuna quantitativa é preenchida por Şakar e Emekci~\cite{sakar2025rag}, que conduzem uma busca em grade de 23.625 iterações comparando seis métodos (Stuff, Refine, Map Reduce, Map Re-rank, Query Step-Down e Reciprocal RAG) e filtros de compressão contextual ao longo de múltiplos bancos vetoriais, modelos de embedding e LLMs. Os autores demonstram que não existe arquitetura universalmente ótima: enquanto o Reciprocal RAG atinge a maior similaridade, o método Stuff é 38,9\% mais rápido e consome 70,5\% menos tokens, e a compressão contextual reduz o uso de recursos ao custo de degradar a similaridade quando seus limiares não são calibrados por domínio. Esse trabalho constitui a principal âncora teórica e metodológica da presente pesquisa.

\subsection{Estratégias avançadas de recuperação}

Para além da simples injeção dos fragmentos recuperados, parte expressiva da literatura busca refinar a etapa de recuperação. Ma et al.~\cite{ma2023query} propõem o paradigma \textit{Rewrite-Retrieve-Read}, no qual um reescritor treinado por aprendizado por reforço adapta a consulta do usuário antes da busca, mitigando a lacuna entre a pergunta e o conhecimento indexado. Em direção complementar, Rackauckas~\cite{rackauckas2024ragfusion} combina a geração de múltiplas consultas com a fusão recíproca de postos (\textit{reciprocal rank fusion}) no método RAG-Fusion, ampliando a abrangência das respostas, embora ao risco de desvio temático quando as consultas geradas se distanciam da intenção original. Essas abordagens elevam a qualidade da recuperação, mas introduzem chamadas adicionais ao LLM, reforçando a tensão entre acurácia e custo que este trabalho pretende mensurar.

\subsection{Contexto longo, estruturação documental e fragmentação}

A forma como o conteúdo recuperado é apresentado ao modelo também impacta o desempenho. Liu et al.~\cite{liu2023lost} demonstram empiricamente o fenômeno \textit{lost in the middle}: o desempenho dos LLMs degrada-se quando a informação relevante se encontra no meio de contextos longos, o que questiona a estratégia de simplesmente acumular fragmentos. Como resposta, Nair et al.~\cite{nair2023drilling} exploram a estrutura de discurso de documentos longos para construir representações condensadas, preservando 99,6\% do desempenho do melhor método zero-shot enquanto processam apenas 26\% dos tokens. No nível mais elementar, Schwaber-Cohen e Patel~\cite{schwaber2025chunking} discutem estratégias de fragmentação (\textit{chunking}) e seu efeito sobre a relevância da recuperação. Esses estudos reforçam que a eficiência em tokens --- objeto deste trabalho --- depende não apenas do método de injeção, mas também da estruturação prévia do contexto.

\subsection{Infraestrutura de armazenamento e recuperação vetorial}

O desempenho de um pipeline RAG é igualmente condicionado pela camada de armazenamento. O levantamento abrangente de Ma et al.~\cite{ma2024vectorsurvey} sistematiza técnicas de armazenamento e recuperação em bancos de dados vetoriais, bem como seus desafios de escalabilidade, fornecendo a base conceitual para a comparação entre soluções como ChromaDB, FAISS e Pinecone prevista nos objetivos deste estudo. A escolha desses componentes, como apontado por Şakar e Emekci~\cite{sakar2025rag}, tem impacto direto sobre o consumo de CPU e a latência, conectando a infraestrutura de recuperação às métricas de hardware aqui investigadas.

\subsection{Aplicações de domínio específico}

Por fim, a literatura evidencia o potencial do RAG em domínios verticais de alta complexidade. Cui et al.~\cite{cui2023chatlaw} apresentam o Chatlaw, assistente jurídico multiagente que combina recuperação de conhecimento legal com uma arquitetura de mistura de especialistas para reduzir alucinações em um domínio sensível. No setor financeiro, Gupta~\cite{gupta2023gptinvestar} utiliza LLMs para analisar relatórios anuais extensos e derivar estratégias de investimento, ilustrando o desafio de processar documentos longos e ricos em jargão. Esses casos confirmam a relevância prática da otimização de RAG e dialogam com os domínios financeiro e técnico empregados como base empírica nesta pesquisa.

\subsection{Lacunas identificadas}

A análise conjunta revela que a literatura avançou significativamente na acurácia~\cite{ma2023query, rackauckas2024ragfusion}, na confiabilidade~\cite{andriopoulos2023augmenting, barnett2024seven} e na compreensão das limitações de contexto~\cite{liu2023lost, nair2023drilling}. Entretanto, a correlação sistemática entre a escolha do método RAG e o impacto físico sobre o hardware --- consumo de CPU, memória RAM e tempo de resposta --- permanece pouco explorada, sendo tratada de forma agregada apenas por Şakar e Emekci~\cite{sakar2025rag} em escala computacional incompatível com ambientes restritos. É justamente nessa lacuna --- mapear de forma reprodutível os trade-offs de eficiência de hardware entre os métodos de RAG --- que este trabalho se insere.
